\documentclass[10pt]{article}
\usepackage[letterpaper,margin=0.42in]{geometry}
\usepackage[T1]{fontenc}
\usepackage{lmodern}
\usepackage{multicol}
\usepackage{enumitem}
\usepackage{xcolor}
\usepackage{titlesec}
\usepackage{microtype}
\usepackage{framed}

\pagestyle{empty}
\setlength{\parindent}{0pt}
\setlength{\columnsep}{0.25in}
\setlist[itemize]{leftmargin=*,nosep}
\setlist[enumerate]{leftmargin=*,nosep}
\titleformat{\section}{\large\bfseries\color{blue!60!black}}{}{0pt}{}[\hrule]
\titlespacing*{\section}{0pt}{4pt}{2pt}
\titleformat{\subsection}{\normalsize\bfseries}{}{0pt}{}
\titlespacing*{\subsection}{0pt}{3pt}{1pt}
\newcommand{\cmd}[1]{\texttt{\small #1}}
\newcommand{\code}[1]{\texttt{#1}}

\begin{document}
\begin{center}
  {\Large\bfseries ICPC NENA Quick Reference}\\[-2pt]
  {\small Bash, C++, Java, graph theory, and contest workflow}
\end{center}

\begin{multicols}{2}

\section{Bash basics}
\subsection{Navigation and files}
\begin{itemize}
  \item \cmd{pwd} --- print the current directory
  \item \cmd{ls -la}, \cmd{cd path}, \cmd{cd ..}, \cmd{cd -}
  \item \cmd{mkdir -p dir}, \cmd{touch file}
  \item \cmd{cp -r src dest}, \cmd{mv old new}
  \item \cmd{rm file}, \cmd{rm -r dir} \textit{(check the path first)}
  \item \cmd{cat file}, \cmd{less file}, \cmd{head -n 20 file}
\end{itemize}

\subsection{Search and inspect}
\begin{itemize}
  \item \cmd{grep -n 'pattern' file}
  \item \cmd{grep -Rni 'pattern' .}
  \item \cmd{find . -type f -name '*.cpp'}
  \item \cmd{wc -l file}, \cmd{sort file}, \cmd{uniq -c file}
  \item \cmd{diff -u old new}
\end{itemize}

\subsection{Compile and run}
\textbf{Set up simple contest commands} in \code{\textasciitilde/.bashrc}:
\scriptsize
\begin{verbatim}
alias compilejava='javac Main.java'
alias runjava='java Main'
alias compileg++='g++ -std=c++17 -O2 -Wall -Wextra -o main'
alias runcpp='./main'
\end{verbatim}
\normalsize
\textbf{Load them} with \cmd{source \textasciitilde/.bashrc}, then use:
\begin{itemize}
  \item Java: \cmd{compilejava}; then \cmd{runjava < input.txt}.
  \item C++: \cmd{compileg++ main.cpp}; then \cmd{runcpp < input.txt}.
  \item Save output: \cmd{runjava < input.txt > output.txt}.
  \item Save output: \cmd{runcpp < input.txt > output.txt}.
  \item Time a run: \cmd{time runcpp < input.txt}.
  \item \cmd{chmod +x script.sh}; \cmd{./script.sh}
\end{itemize}

\subsection{Pipes and shell operators}
\begin{itemize}
  \item \cmd{cmd1 | cmd2} --- send output to the next command
  \item \cmd{cmd > file} --- overwrite; \cmd{cmd >> file} --- append
  \item \cmd{cmd1 \&\& cmd2} --- run the second only on success
  \item \cmd{cmd1 || cmd2} --- run the second only on failure
  \item \cmd{cmd1; cmd2} --- run both regardless of status
  \item \cmd{Ctrl-C} stop; \cmd{Ctrl-D} end input; \cmd{Ctrl-L} clear
\end{itemize}

\subsection{Variables and loops}
\begin{itemize}
  \item \cmd{X=42}; use with \cmd{\$X}; export with \cmd{export X}
  \item \cmd{for f in *.cpp; do echo "\$f"; done}
  \item \cmd{while read -r line; do echo "\$line"; done < file}
  \item \cmd{echo "text"}; \cmd{printf '\%s\textbackslash n' "\$value"}
\end{itemize}

\section{Contest workflow}
\begin{enumerate}
  \item Confirm the working directory with \cmd{pwd}.
  \item Compile with warnings enabled; fix every warning you understand.
  \item Test edge cases: empty/minimum, maximum, duplicates, overflow.
  \item Compare output using \cmd{diff -u expected output}.
  \item Keep source, input, and output files in a clearly named directory.
\end{enumerate}

\section{Useful environment commands}
\begin{itemize}
  \item \cmd{which g++}, \cmd{g++ --version}, \cmd{python3 --version}
  \item \cmd{echo \$PATH}, \cmd{env}
  \item \cmd{ps}, \cmd{kill PID}
  \item \cmd{man command}, \cmd{command --help}
\end{itemize}

\section{C++ compile reminders}
\begin{itemize}
  \item Prefer \code{long long} for values up to about $9\times10^{18}$.
  \item Use \code{ios::sync\_with\_stdio(false);}\\
        \code{cin.tie(nullptr);} for fast I/O.
  \item Compile locally with the same standard and flags used for submission.
\end{itemize}

\section{More Bash tools}
\subsection{Redirection and files}
\begin{itemize}
  \item \cmd{tee output.txt} displays output and saves a copy.
  \item \cmd{sort -n file}; \cmd{sort -rn file}; \cmd{uniq -c}.
  \item \cmd{du -sh .}; \cmd{df -h} check disk usage and free space.
  \item \cmd{realpath file}; \cmd{basename path}; \cmd{dirname path}.
\end{itemize}

\subsection{Permissions and processes}
\begin{itemize}
  \item \cmd{chmod +x script.sh} makes a script executable.
  \item \cmd{ps aux | grep main} finds a running program.
  \item \cmd{kill PID} stops a process; use \cmd{Ctrl-C} first.
  \item \cmd{history}; repeat the last command with \cmd{!!}.
\end{itemize}

\subsection{Repeatable testing}
\begin{itemize}
  \item \cmd{for i in \{1..100\}; do ./main < in; done}
  \item \cmd{diff -w expected output} ignores whitespace differences.
  \item \cmd{timeout 2s ./main < input.txt} catches very slow programs.
  \item \cmd{time -p ./main < input.txt} reports runtime.
\end{itemize}

\end{multicols}
\newpage
\begin{center}
  {\Large\bfseries Java Quick Reference}
\end{center}
\begin{multicols}{2}

\section{Java}
\subsection{Compile and run}
\begin{itemize}
  \item File must be named \code{Main.java}; class must be \code{Main}.
  \item \cmd{complilejava Main.java}; \cmd{runjava Main}
  \item Run with redirected input and output:
    \cmd{java Main < input.txt > output.txt}
\end{itemize}

\end{multicols}

\subsection{Fast input template}
\begin{itemize}
  \item \code{BufferedReader} and \code{StringTokenizer} avoid the overhead of \code{Scanner}.
  \item Buffer output with \code{StringBuilder} or \code{PrintWriter}.
\end{itemize}
\begin{framed}
\noindent\textbf{FastReader.java}\par\smallskip
\scriptsize
\setlength{\baselineskip}{7.2pt}
\begin{verbatim}
// FastReader class for efficient input
static class FastReader {
  
    // BufferedReader to read input
    BufferedReader b;
  
    // StringTokenizer to tokenize input
    StringTokenizer s; 

    // Constructor to initialize BufferedReader
    public FastReader() {
        b = new BufferedReader(new InputStreamReader(System.in));
    }

    // Method to read the next token as a string
    String next() {
        while (s == null || !s.hasMoreElements()) {
            try {
                s = new StringTokenizer(b.readLine());
            } catch (IOException e) {
                e.printStackTrace(); 
            }
        }
        return s.nextToken();
    }

    // Method to read the next token as an integer
    int nextInt() { 
        return Integer.parseInt(next()); 
    }

    // Method to read the next token as a long
    long nextLong() { 
        return Long.parseLong(next()); 
    }

    // Method to read the next token as a double
    double nextDouble() { 
        return Double.parseDouble(next()); 
    }

    // Method to read the next line as a string
    String nextLine() {
        String str = "";
        try {
            if (s.hasMoreTokens()) {
                str = s.nextToken("\n");
            } else {
                str = b.readLine();
            }
        } catch (IOException e) {
            e.printStackTrace(); 
        }
        return str;
  }
}
\end{verbatim}
\end{framed}

\subsection{Fast output}
\begin{itemize}
  \item Build all output, then print once:
\end{itemize}
\begin{framed}
\footnotesize
\begin{verbatim}
StringBuilder out = new StringBuilder();
out.append(answer).append('\n');
out.append(value).append(' ').append(other).append('\n');
System.out.print(out);
\end{verbatim}
\end{framed}
\begin{itemize}
  \item Or use \code{PrintWriter out = new PrintWriter(System.out);}
  \item Finish with \code{out.println(answer); out.flush();}.
\end{itemize}

\begin{multicols}{2}
\subsection{Common utilities and pitfalls}
\begin{itemize}
  \item Sort arrays: \code{Arrays.sort(a)}.
    Sort lists: \code{Collections.sort(list)}.
  \item Lists: \code{ArrayList<Integer>}; maps:
    \code{HashMap<K,V>}; sets: \code{HashSet<T>}.
  \item Priority queue: \code{PriorityQueue<Integer> pq}
  \item Use \code{long}; integer division truncates, so cast before division.
  \item Max-heap: \code{new PriorityQueue<>}\\
    \code{(Collections.reverseOrder())}.
  \item Use \code{Math.floorMod(a, m)} for a nonnegative remainder.
\end{itemize}

\end{multicols}
\begin{center}
  {\Large\bfseries Graph Theory Algorithms}
\end{center}
\begin{multicols}{2}

\section{Graph representation}
\begin{itemize}
  \item Adjacency list: \code{List<Integer>[] g}.
  \item Weighted edge: \code{class Edge \{ int to, w; \}}.
  \item Undirected edge $(u,v)$: add both \code{u -> v} and \code{v -> u}.
  \item Store vertices as $0,\ldots,n-1$ and use a \code{visited} array.
\end{itemize}

\section{Traversal}
\begin{itemize}
  \item \textbf{BFS:} shortest paths in an unweighted graph.
    Use a queue; time $O(V+E)$.
  \item \textbf{DFS:} connectivity, components, cycle checks, and ordering.
    Time $O(V+E)$.
  \item For a bipartite check, color each newly discovered vertex with
    the opposite color; a same-color edge is a contradiction.
\end{itemize}

\section{Shortest paths}
\begin{itemize}
  \item \textbf{Dijkstra:} nonnegative weights; priority queue;
    $O((V+E)\log V)$.
  \item \textbf{0--1 BFS:} weights only $0$ or $1$; deque; $O(V+E)$.
  \item \textbf{Bellman--Ford:} handles negative edges and detects negative
    cycles; $O(VE)$.
  \item \textbf{Floyd--Warshall:} all pairs; $O(V^3)$ and $O(V^2)$ memory.
  \item Use a large \code{INF}; never add to \code{INF} without checking.
\end{itemize}

\section{Minimum spanning tree}
\begin{itemize}
  \item \textbf{Kruskal:} sort edges, add if endpoints are in different
    DSU sets; $O(E\log E)$.
  \item \textbf{Prim:} grow from one vertex with a min-heap;
    $O((V+E)\log V)$ with adjacency lists.
  \item DSU operations: path compression and union by size/rank are nearly
    constant amortized time.
\end{itemize}

\section{Directed graphs}
\begin{itemize}
  \item \textbf{Topological sort:} DAG ordering. Kahn's algorithm repeatedly
    removes indegree-zero vertices; $O(V+E)$.
  \item If fewer than $V$ vertices are removed, the graph contains a cycle.
  \item \textbf{SCCs:} Kosaraju uses DFS on the graph and its reverse;
    Tarjan finds SCCs in $O(V+E)$.
\end{itemize}

\section{Useful reminders}
\begin{itemize}
  \item A tree with $V$ vertices has exactly $V-1$ edges.
  \item Connected undirected graph: DFS/BFS reaches every vertex.
  \item A connected graph has an Euler circuit iff every vertex has even
    degree; it has an Euler path iff exactly zero or two have odd degree.
  \item For a tree, the diameter can be found with two BFS/DFS runs.
\end{itemize}

\end{multicols}

\vfill
\hrule
\small
\textbf{Printing note:} This sheet is intentionally compact and uses only standard
LaTeX packages. Confirm the current ICPC NENA rules for permitted reference
materials, page count, paper size, and annotations before bringing it to a contest.

\end{document}
